\chapter{Machine Learning con scikit-learn}
\label{chap:machine-learning}
{\huge$\bigstar$}{\huge$\bigstar$}{\huge$\bigstar$}{\huge$\bigstar$} Experto

El ecosistema Python ofrece las bibliotecas de machine learning más completas
del mundo. Con LisPy, puedes usar scikit-learn, pandas, y NumPy directamente
desde Lisp, combinando la expresividad de Lisp con el poder computacional
de Python.

\section{Clasificación con Random Forest}

El clasificador Random Forest es uno de los algoritmos más versátiles.
Funciona construyendo múltiples árboles de decisión durante el entrenamiento:

\begin{lispcode}
(py:run "from sklearn.ensemble import RandomForestClassifier
from sklearn.datasets import load_iris
from sklearn.model_selection import train_test_split
from sklearn.metrics import accuracy_score, classification_report
iris = load_iris()
X_train, X_test, y_train, y_test = train_test_split(
    iris.data, iris.target, test_size=0.2, random_state=42)
clf = RandomForestClassifier(n_estimators=100)
clf.fit(X_train, y_train); y_pred = clf.predict(X_test)
print('Accuracy:', accuracy_score(y_test, y_pred))
print(classification_report(y_test, y_pred, target_names=iris.target_names))")
\end{lispcode}

La salida incluye la precisión (\texttt{accuracy}) y un reporte detallado
con precisión por clase, exhaustividad (\texttt{recall}), y $F_1$-score.

\section{Pipelines de Preprocesamiento}

Los pipelines permiten encadenar transformaciones y el clasificador final:

\begin{lispcode}
(py:run "from sklearn.pipeline import Pipeline
from sklearn.preprocessing import StandardScaler
from sklearn.decomposition import PCA
from sklearn.svm import SVC
from sklearn.datasets import load_digits
from sklearn.model_selection import train_test_split
digits = load_digits()
X_train, X_test, y_train, y_test = train_test_split(
    digits.data, digits.target, test_size=0.2, random_state=42)
pipeline = Pipeline([('scaler', StandardScaler()),
                     ('pca', PCA(n_components=30)),
                     ('svc', SVC(kernel='rbf'))])
pipeline.fit(X_train, y_train)
print('Pipeline score:', pipeline.score(X_test, y_test))")
\end{lispcode}

\begin{center}
\begin{tikzpicture}[node distance=2cm, auto]
  \tikzstyle{block} = [rectangle, draw, fill=lispblue!10,
    text width=3cm, text centered, rounded corners, minimum height=1cm]
  \tikzstyle{arrow} = [->, >=stealth]

  \node [block] (data) {Datos crudos};
  \node [block, right of=data, node distance=3.5cm] (scaler)
    {Standard\-Scaler};
  \node [block, right of=scaler, node distance=3.5cm] (pca)
    {PCA (30 comp.)};
  \node [block, right of=pca, node distance=3.5cm] (svc)
    {SVC (RBF)};
  \node [block, right of=svc, node distance=3.5cm] (pred)
    {Predicción};

  \draw [arrow] (data) -- (scaler);
  \draw [arrow] (scaler) -- (pca);
  \draw [arrow] (pca) -- (svc);
  \draw [arrow] (svc) -- (pred);
\end{tikzpicture}
\end{center}

El pipeline asegura que la estandarización y la reducción de dimensionalidad
se apliquen de manera consistente en entrenamiento y prueba, evitando
fugas de datos (\texttt{data leakage}).

\section{Validación Cruzada}

Para una evaluación más robusta, usamos validación cruzada con
\texttt{cross\_val\_score}:

\begin{lispcode}
(py:run "from sklearn.model_selection import cross_val_score
from sklearn.ensemble import RandomForestClassifier
from sklearn.datasets import load_iris
iris = load_iris()
clf = RandomForestClassifier(n_estimators=50)
scores = cross_val_score(clf, iris.data, iris.target, cv=5)
print('Scores:', scores)
print('Mean:', scores.mean(), '+/-', scores.std() * 2)")
\end{lispcode}

\section{Clustering K-Means}

El clustering no supervisado permite descubrir grupos en datos no etiquetados:

\begin{lispcode}
(py:run "from sklearn.cluster import KMeans
from sklearn.datasets import make_blobs
import matplotlib.pyplot as plt
X, y_true = make_blobs(n_samples=300, centers=4,
                       cluster_std=0.6, random_state=42)
kmeans = KMeans(n_clusters=4, random_state=42)
y_kmeans = kmeans.fit_predict(X)
plt.scatter(X[:,0], X[:,1], c=y_kmeans, cmap='viridis')
plt.scatter(kmeans.cluster_centers_[:,0],
            kmeans.cluster_centers_[:,1],
            marker='x', s=200, linewidths=3, color='red')
plt.title('K-Means Clustering (k=4)')
plt.savefig('kmeans.png')")
\end{lispcode}

\begin{center}
\begin{tikzpicture}
\begin{axis}[
  width=0.85\textwidth,
  height=7cm,
  xlabel={Característica 1},
  ylabel={Característica 2},
  title={K-Means: 4 clusters},
  grid=both,
  only marks,
  mark size=1.5,
  scatter,
  point meta=y,
]
\addplot[scatter, color=lispblue] table {
0.5 0.5 0
0.6 0.7 0
1.2 1.1 0
4.5 4.3 1
4.7 4.0 1
5.0 4.5 1
8.1 0.2 2
7.9 0.5 2
8.5 0.0 2
2.0 7.5 3
2.3 7.8 3
1.8 8.0 3
};
\end{axis}
\end{tikzpicture}
\end{center}

\section{Regresión Lineal}

Para problemas de regresión, scikit-learn ofrece modelos lineales y
regularizados:

\begin{lispcode}
(py:run "from sklearn.linear_model import LinearRegression
from sklearn.datasets import make_regression
from sklearn.metrics import mean_squared_error, r2_score
import numpy as np
X, y = make_regression(n_samples=100, n_features=1, noise=10,
                       random_state=42)
X_train, X_test = X[:80], X[80:]
y_train, y_test = y[:80], y[80:]
lr = LinearRegression(); lr.fit(X_train, y_train)
y_pred = lr.predict(X_test)
print('MSE:', mean_squared_error(y_test, y_pred))
print('R2:', r2_score(y_test, y_pred))
print('Coef:', lr.coef_[0], 'Intercept:', lr.intercept_)")
\end{lispcode}

\section{Búsqueda de Hiperparámetros}

Optimización sistemática con GridSearchCV:

\begin{lispcode}
(py:run "from sklearn.model_selection import GridSearchCV
from sklearn.ensemble import RandomForestClassifier
from sklearn.datasets import load_iris
iris = load_iris()
param_grid = {'n_estimators': [50, 100, 200],
              'max_depth': [3, 5, None]}
grid = GridSearchCV(RandomForestClassifier(random_state=42),
                    param_grid, cv=5, scoring='accuracy')
grid.fit(iris.data, iris.target)
print('Best params:', grid.best_params_)
print('Best score:', grid.best_score_)")
\end{lispcode}

\section{Datos Tabulares con pandas}

Los datos del mundo real rara vez llegan en el formato exacto que necesita
un modelo. pandas proporciona las herramientas para cargar, limpiar y
transformar datos:

\begin{lispcode}
(py:run "import pandas as pd; from sklearn.ensemble import RandomForestRegressor
from sklearn.model_selection import train_test_split
from sklearn.metrics import mean_absolute_error
df = pd.read_csv('datos.csv')
X = df.drop('target', axis=1); y = df['target']
X_train, X_test, y_train, y_test = train_test_split(X, y, test_size=0.2)
model = RandomForestRegressor(n_estimators=100)
model.fit(X_train, y_train)
preds = model.predict(X_test)
print('MAE:', mean_absolute_error(y_test, preds))")
\end{lispcode}

\section{Exportación de Modelos}

Los modelos entrenados se pueden serializar con pickle para uso posterior:

\begin{lispcode}
(py:run "import pickle; import pandas as pd
from sklearn.ensemble import RandomForestClassifier
from sklearn.datasets import load_iris
iris = load_iris(); clf = RandomForestClassifier()
clf.fit(iris.data, iris.target)
with open('model.pkl', 'wb') as f: pickle.dump(clf, f)
with open('model.pkl', 'rb') as f:
    loaded = pickle.load(f)
print('Loaded model score:', loaded.score(iris.data, iris.target))")
\end{lispcode}

\section{Resumen de Flujo de Trabajo}

El flujo típico de un proyecto de machine learning con LisPy sigue estos
pasos:

\begin{enumerate}
\item \textbf{Carga de datos}: pandas (CSV, JSON, SQL, Excel) o NumPy
\item \textbf{Preprocesamiento}: StandardScaler, LabelEncoder, PCA, selección de características
\item \textbf{División}: train\_test\_split con estratificación
\item \textbf{Modelado}: RandomForest, SVC, KMeans, LinearRegression, etc.
\item \textbf{Evaluación}: cross\_val\_score, métricas (accuracy, MSE, $F_1$)
\item \textbf{Optimización}: GridSearchCV, RandomizedSearchCV
\item \textbf{Despliegue}: pickle, ONNX, o integración en pipeline Lisp
\end{enumerate}

Cada uno de estos pasos puede ser orquestado desde Lisp, permitiendo
generar pipelines de ML de manera programática con macros y metaprogramación.

\endinput
